\bookchapter{N100: Grammar and Exact
Work}\label{viii-ch01-n100-grammar-and-exact-work}

The first computational question is exact and finite: what object is
specified by the admitted source grammar, and how can its complete state
sum be evaluated? This chapter keeps the object fixed while following
two execution routes. The source census establishes the object; the
factorization explains why its formal cardinality need not equal the
amount of local enumeration.

The account below follows the chapter source as retained at the fixed 8
September 2026 SAMA revision. Engine names, installed interfaces and
campaign status are historical source statements. They do not describe a
freshly inspected October deployment. Later source amendments are
distinguished where available; no later engine or campaign is inferred.

\subsection{Overview}\label{viii-ch01-overview}

N100 is an exact frustrated-spin computation whose favorable
computational structure comes from its source grammar. The mathematical
object is the unchanged 100-spin, 260-edge Hamiltonian with

\[
2^{100}=1{,}267{,}650{,}600{,}228{,}229{,}401{,}496{,}703{,}205{,}376
\]

binary assignments. The first complete H14F route returns the full exact
density-of-states tensor in 3.828137395 seconds. A later structure-first
route detects eight conditional components, executes 165,920 local
assignments and convolves them to the coefficientwise-identical tensor
in 0.500951441 seconds. The reduced workload is an exact evaluation
strategy; it does not reduce the state space, change the Hamiltonian,
sample the assignments or request only a ground state.

The 100-row census is itself source-bound. The visible QP lane contains
105 records. Exactly one of the 21 complete five-role packet
exclusions---the \texttt{p=9,g=0} packet---produces 100 independently
counted rows with exact native inventory 16,200 and mean 162. The five
excluded-count W9 witnesses remain visible as derived \texttt{X1+W8}
receipts. The A field remains present as non-row infrastructure, and X1
remains non-row/non-matter. Those are three different boundaries.

The central computational finding is not merely the observed wall time.
The source grammar contains exact decomposition certificates, and the
compiler can use them while preserving the complete result. Its workload
also returns the ledger quantity L162 through an independent arithmetic
identity. The N105, N81 and N72 controls show exactly which parts depend
on the source grammar and which do not.

\subsection{Opening question and conceptual
picture}\label{viii-ch01-opening-question-and-conceptual-picture}

The governing question is:

\begin{quote}
Can a source-derived 100-row grammar expose enough exact structure to
compute the complete \texttt{2\^{}100} density-of-states tensor without
visiting every assignment one by one, while returning exactly the same
mathematical output?
\end{quote}

The answer is represented by two equivalent exact routes:

\begin{itemize}
\tightlist
\item
  source-derived N100 Hamiltonian
\item
  complete H14F variable-elimination route
\item
  exact retained-role DOS tensor
\item
  same unchanged Hamiltonian
\item
  detect exact conditional components
\item
  compute each local exact tensor
\item
  exact tensor/polynomial convolution
\item
  coefficientwise-identical retained-role DOS tensor
\end{itemize}

The grammar does not replace the calculation. It tells the exact solver
how the complete sum can be factorized.

\subsection{Definitions and cardinality
firewall}\label{viii-ch01-definitions-and-cardinality-firewall}

\subsubsection{Exact frustrated-spin density of
states}\label{viii-ch01-exact-frustrated-spin-density-of-states}

For binary spins

\[
\sigma_i\in\{-1,+1\},\qquad i=1,\ldots,N,
\]

and a declared weighted graph, the exact Hamiltonian has the form

\[
H(\sigma)
=\sum_{(i,j)\in E}J_{ij}\sigma_i\sigma_j
+\sum_i h_i\sigma_i.
\]

The density of states is the integer function

\[
g(E)=\#\{\sigma\in\{-1,+1\}^N:H(\sigma)=E\}.
\]

Completeness requires

\[
\sum_Eg(E)=2^N.
\]

For N100, \texttt{N=100}. Frustration changes the organization of energy
and the ground sector; it does not change the formal assignment count.

\subsubsection{Four quantities that must not be
merged}\label{viii-ch01-four-quantities-that-must-not-be-merged}

\begin{itemize}
\tightlist
\item
  \textbf{N100 source rows.} Exact value: 100; Type: independent spins
  of the declared Hamiltonian.
\item
  \textbf{complete configuration space.} Exact value: \texttt{2\^{}100};
  Type: every formal binary assignment counted by the returned tensor.
\item
  \textbf{structure-first work.} Exact value: 165,920; Type: local
  assignments evaluated before exact convolution.
\item
  \textbf{L162 native ledger mean.} Exact value: 162; Type:
  source-ledger/accounting identity, not a coupling or state count.
\end{itemize}

The statement ``exact \texttt{2\^{}100}'' refers to the tensor's
complete count. It does not say the runner iterated through
\texttt{2\^{}100} bit strings. Conversely, the 165,920 workload is not a
claim that the physical or mathematical problem had only 165,920 states.

\subsubsection{The row-census exclusions are separately
typed}\label{viii-ch01-the-row-census-exclusions-are-separately-typed}

The N100 census is produced after one historical five-row packet is
excluded from independent content. Preserve these three boundaries:

\begin{itemize}
\tightlist
\item
  p=9,g=0 packet: five derived W9 witnesses, visible but non-counted
\item
  A-field row: excluded because A is active non-row infrastructure
\item
  X1: response channel, not a physical/source row
\end{itemize}

Similar arithmetic around the 163-to-162 ledger does not make the three
objects equivalent.

\subsubsection{Source values are not Hamiltonian
couplings}\label{viii-ch01-source-values-are-not-hamiltonian-couplings}

The QP rows supply typed identity, topology and role information through
a separately declared compiler. Their native numerical values are not
inserted as \texttt{J\_ij} or \texttt{h\_i} merely because they occur on
the same record. The tested Hamiltonian is the source-derived 100-spin,
260-edge graph fixed before the structure-first route.

\subsection{Deriving the 100-row source
census}\label{viii-ch01-deriving-the-100-row-source-census}

\subsubsection{Begin with the full visible
lane}\label{viii-ch01-begin-with-the-full-visible-lane}

CR120R begins from 105 visible source rows arranged into 21 complete
five-role address packets. The exact source totals are

\begin{itemize}
\tightlist
\item
  source rows 105
\item
  complete five-role packets 21
\item
  address-value sum 2889
\item
  full native inventory 130005/8 = 16250.625
\end{itemize}

Every source row remains present in the result artifact. The operation
is a counted/not-counted overlay, not deletion.

\subsubsection{Test every one-packet
exclusion}\label{viii-ch01-test-every-one-packet-exclusion}

Each of the 21 five-role packets is removed in turn from
independent-content accounting. A valid candidate must leave 100 rows
and close the L162 mean:

\[
\frac{M_{native,remaining}}{100}=162.
\]

Exactly one candidate does so:

\[
(p,g)=(9,0).
\]

The resulting payload is

\begin{itemize}
\tightlist
\item
  excluded typed packet p=9,g=0, five rows
\item
  payload rows 100
\item
  payload native inventory 16200
\item
  payload mean 162=L
\item
  single-packet L162 hits 1
\end{itemize}

Thus

\[
\frac{16{,}200}{100}=162.
\]

\subsubsection{Verify the derived-parent
relation}\label{viii-ch01-verify-the-derived-parent-relation}

Each same-role \texttt{p=9,g=0} record closes exactly as the sum of its
\texttt{p=8,g=0} and \texttt{p=1,g=0} parents in \texttt{M\_native}.
This is the source-typed W9 route

\[
8+1\to9
\]

expressed on five records. The type controls prevent the scalar 9 from
collapsing \texttt{W9\_CLOSURE\_WITNESS}, \texttt{C9\_CARRIER\_ATOM} and
\texttt{P9\_BIGRADE\_PARTITION} into one object. Weak-vector and
hidden-source p9 support rows are also barred from the selector, and the
identity is not extended to g1, g2, composite routes or qA.

\subsubsection{Reconcile the payload}\label{viii-ch01-reconcile-the-payload}

The derived 100 payload IDs and values reconcile one-for-one with the
current \nolinkurl{Updated\ Particle\ Rows.xlsx} extraction recorded by
CR120R. This closes the source census used by the N100 computation while
retaining the five witnesses as provenance.

\subsection{Constructing the exact N100
problem}\label{viii-ch01-constructing-the-exact-n100-problem}

\subsubsection{Freeze the graph and requested
output}\label{viii-ch01-freeze-the-graph-and-requested-output}

The declared problem contains:

\begin{itemize}
\tightlist
\item
  100 independent binary spins;
\item
  260 source-derived edges;
\item
  induced width 5 under the complete H14F route;
\item
  a retained-role conditional DOS tensor;
\item
  a 16-row scalar closure; and
\item
  exact integer coefficients reconstructed across four CRT primes.
\end{itemize}

The requested output is the complete tensor, not only the minimum energy
or a small collection of moments. The graph, source roles and tensor
schema are fixed before either timing result.

\subsubsection{Why exact routes can be smaller than the state
count}\label{viii-ch01-why-exact-routes-can-be-smaller-than-the-state-count}

The complete sum

\[
g(E)=\sum_{\sigma\in\{-1,+1\}^{100}}
\mathbf 1\{H(\sigma)=E\}
\]

can be rearranged without approximation. Variable elimination sums
internal spins while retaining exact boundary spectra. Disconnected
conditional components can be computed independently and convolved.
Modular NTT lanes can multiply polynomials exactly, and CRT reconstructs
the integer coefficients once the modulus product exceeds the
coefficient range.

All such routes remain exact if they reproduce every coefficient and
close the complete state count.

\subsection{Complete H14F route from execution to
closure}\label{viii-ch01-complete-h14f-route-from-execution-to-closure}

\subsubsection{Execute the declared
Hamiltonian}\label{viii-ch01-execute-the-declared-hamiltonian}

The first route applies H14F variable elimination to the full
source-derived graph. It produces:

\begin{itemize}
\tightlist
\item
  \textbf{spins.} Exact result: 100.
\item
  \textbf{source edges.} Exact result: 260.
\item
  \textbf{represented configurations.} Exact result: \texttt{2\^{}100}.
\item
  \textbf{occupied energy bins.} Exact result: 177.
\item
  \textbf{ground sector.} Exact result: \texttt{(-260,256)}.
\item
  \textbf{distinct port spectra.} Exact result: 3.
\item
  \textbf{singleton root tasks.} Exact result: 2,048.
\item
  \textbf{CRT primes.} Exact result: 4.
\item
  \textbf{observed runner wall.} Exact result: 3.828137395 s.
\end{itemize}

The ordered pair \texttt{(-260,256)} records the ground energy and its
degeneracy. The calculation retains the other 176 occupied bins as well.

\subsubsection{Close the port marginals}\label{viii-ch01-close-the-port-marginals}

Four boundary bits produce 16 retained port rows. Each row sums to

\[
2^{96},
\]

because fixing the four retained bits leaves 96 independent spins.
Summing all 16 rows gives

\[
16\cdot2^{96}=2^4\cdot2^{96}=2^{100}.
\]

This is both a marginal check and a full-state closure.

\subsubsection{Close source symmetries}\label{viii-ch01-close-source-symmetries}

All 24 source-role permutations reproduce the required tensor relation.
The global spin complement also closes. These checks distinguish a
complete source-role tensor from an accidental scalar DOS total.

\subsection{Structure-first route from components to the same
tensor}\label{viii-ch01-structure-first-route-from-components-to-the-same-tensor}

\subsubsection{Detect exact conditional
components}\label{viii-ch01-detect-exact-conditional-components}

The source grammar exposes eight components with sizes

\[
15,15,15,15,15,10,10,5.
\]

They are conditional components: any shared retained port variables
remain explicit, while the interior variables of each component can be
summed locally.

\subsubsection{Compute local exact
tensors}\label{viii-ch01-compute-local-exact-tensors}

For component \texttt{c} and retained boundary assignment \texttt{b},
define its exact energy polynomial

\[
P_{c,b}(x)=\sum_E g_{c,b}(E)x^E.
\]

Exhaustive local work for the eight interiors is

\[
W_{100}
=5\cdot2^{15}+2\cdot2^{10}+2^5.
\]

Evaluate each term:

\[
5\cdot2^{15}=5\cdot32{,}768=163{,}840,
\]

\[
2\cdot2^{10}=2{,}048,
\qquad2^5=32,
\]

so

\[
W_{100}=163{,}840+2{,}048+32=165{,}920.
\]

\subsubsection{Convolve rather than discard cross-component
sums}\label{viii-ch01-convolve-rather-than-discard-cross-component-sums}

For each retained boundary assignment, multiply the local polynomials:

\[
P_b(x)=\prod_{c=1}^{8}P_{c,b}(x).
\]

Polynomial multiplication performs the full energy convolution: every
choice of one local state from each component contributes to the
coefficient at the sum of their energies. Summing or assembling over the
retained boundary assignments returns the same retained-role tensor as
the complete route.

No state is sampled or dropped. Factorization evaluates the product
structure of the full sum.

\subsubsection{Compare complete outputs}\label{viii-ch01-compare-complete-outputs}

The structure-first wall is 0.500951441 seconds. Relative to the
3.828137395 second H14F route, the observed ratio is approximately

\[
\frac{3.828137395}{0.500951441}\approx7.642.
\]

The load-bearing comparison is stronger than timing: the complete
tensors have the same coefficient hash and match coefficient by
coefficient. The Hamiltonian, state space, energy support and output
schema are unchanged.

\subsection{Grammar-to-work identity, derived line by
line}\label{viii-ch01-grammar-to-work-identity-derived-line-by-line}

The source ledger supplies

\[
N100_{native}=100\cdot162=16{,}200.
\]

The exact component work can be regrouped as

\[
165{,}920=165{,}888+32=162\cdot1{,}024+32.
\]

Therefore

\[
W_{100}=162\cdot2^{10}+2^5.
\]

Subtract the residual and divide by the ten-bit block:

\[
\frac{W_{100}-2^5}{2^{10}}
=\frac{165{,}920-32}{1{,}024}
=\frac{165{,}888}{1{,}024}
=162.
\]

There is a second exact form. Because

\[
72^2=5{,}184=32\cdot162,
\]

we have

\[
W_{100}=2^5(72^2+1)=32(5{,}184+1)=165{,}920.
\]

The residual \texttt{2\^{}5} is the binary assignment space of the
surviving five-role \texttt{p=9,g=1} packet. It is not X1, Theta18 or a
physical closure object. The identity connects two independently sourced
structures---the row ledger and the exact workload---but the value
16,200 is not inserted into the Hamiltonian and does not mechanically
cause the runtime.

\subsection{Source-native controls and
deviations}\label{viii-ch01-source-native-controls-and-deviations}

\subsubsection{N105 restored-packet
control}\label{viii-ch01-n105-restored-packet-control}

Restore the historical \texttt{p=9,g=0} packet. The workload becomes

\[
W_{105}=166{,}912=163\cdot2^{10}.
\]

The difference is

\[
\Delta W=166{,}912-165{,}920=992=2^{10}-2^5.
\]

The restoration also changes the N100/N81 completion-path capacity from
40 to 45 rows. The clean 162 identity is therefore sensitive to the
typed packet boundary; it is not an identity forced by any 100-variable
benchmark.

\subsubsection{N81 octahedral control}\label{viii-ch01-n81-octahedral-control}

The N81 source-native control has component sizes

\[
[11,10,11,10,10,11,10,8]
\]

and exact work

\[
10{,}496=2^7(81+1).
\]

Its octahedral bridge retains 48 cube embeddings in one full-octahedral
orbit, the completion path \texttt{1-8-9-12}, reflected C4 gluing and
49,152 incidences. The N100/N81 localization remains
\texttt{40-\textgreater{}40} under all nine insertions. These are exact
geometry/compiler leads; a causal account of the timing is still a
separate research question.

\subsubsection{N72 dense contrast}\label{viii-ch01-n72-dense-contrast}

N72 remains an independently dense computation surface. It is not edited
to look like N100. This control prevents the source-structure finding
from being rephrased as a universal rule that nominally smaller N must
be faster or that all graphs admit the same component decomposition.

\subsubsection{Unchanged-parent control}\label{viii-ch01-unchanged-parent-control}

The N100 graph and requested tensor were frozen before the
structure-first timing route. The earlier H14F result already supplied
the complete tensor. Coefficientwise equality therefore tests an
optimization of representation on an unchanged parent rather than a
retuned benchmark.

\subsection{Source-native compiler
boundary}\label{viii-ch01-source-native-compiler-boundary}

The complete N100 computation and the later N100-to-F81 semantic
compiler are different tasks. G1 audits the latter and records:

\begin{itemize}
\tightlist
\item
  \textbf{typed N100 QP source keys.} Coverage: 100/100.
\item
  \textbf{Shared72 aggregate connector-incidence eligibility.} Coverage:
  72/100.
\item
  \textbf{CR211 address rows.} Coverage: 162.
\item
  \textbf{CR211 rows with QP elementwise assignment.} Coverage: 0.
\item
  \textbf{F81 semantic sites.} Coverage: 81.
\item
  \textbf{frozen F81 contact classes.} Coverage: 10.
\item
  \textbf{N100 rows with complete elementwise F81 map.} Coverage: 0.
\end{itemize}

The missing key is

\begin{itemize}
\tightlist
\item
  \nolinkurl{QP_CANDIDATE_OR_TEMPLATE_TO_F81_SEMANTIC_SITE_OR_SOURCE_NATIVE_F81_ADDRESS_CLASS}
\end{itemize}

CR211 explicitly leaves address-to-template incidence null, and the F81
target schema has no QP source key. G1 therefore records
\nolinkurl{PARTIAL_DOMAIN__ELEMENTWISE_F81_JOIN_OPEN}; the contact
lane is
\nolinkurl{NOT_EXECUTED__SOURCE_NATIVE_ELEMENTWISE_COMPILER_OPEN}.
No result or residual selects a map.

This boundary does not retract the exact N100 Hamiltonian calculation.
It says that a later semantic join from each source row to an F81
site/address class has not been supplied. Aggregate Shared72 eligibility
cannot substitute for 100 elementwise assignments.

\subsection{Established result and current
boundaries}\label{viii-ch01-established-result-and-current-boundaries}

The source-bound N100 result is:

\begin{enumerate}
\def\labelenumi{\arabic{enumi}.}
\tightlist
\item
  CR120R fixes a 100-row independent-content census after the unique
  \texttt{p=9,g=0} five-record witness packet is excluded from counting.
\item
  The A field remains non-row infrastructure, X1 remains
  non-row/non-matter, and all 105 records remain visible.
\item
  The unchanged Hamiltonian has 100 spins, 260 edges and a complete
  \texttt{2\^{}100} configuration space.
\item
  H14F returns the full exact retained-role DOS tensor with 177 occupied
  bins and ground sector \texttt{(-260,256)}.
\item
  The structure-first route evaluates 165,920 local assignments and
  exactly convolves them to the same coefficient tensor.
\item
  Work reduction is not state-space reduction.
\item
  The exact grammar-to-work identity returns 162 while retaining its
  \texttt{2\^{}5} residual as the p9,g1 binary packet.
\item
  N105, N81 and N72 preserve the source dependence and scope.
\item
  The source-native elementwise N100-to-F81 compiler join remains open,
  and the contact lane was not executed.
\item
  Historical SLC revision labels in the immutable spine do not replace
  the then-live \texttt{SLCQ2-RZ} source-recorded default identity.
\end{enumerate}

The registered atomic records and focused tests assign no authorized
project result classification to this chapter. Exact source verdicts and
timings are preserved without manufacturing one or extending the result
to arbitrary 100-spin graphs.

\subsection{Connections and forward
handoff}\label{viii-ch01-connections-and-forward-handoff}

\begin{itemize}
\tightlist
\item
  \textbf{\texttt{SAMA-D000027}.} Connection: Supplies the finite source
  grammar and row-count distinctions..
\item
  \textbf{\texttt{SAMA-D000033}.} Connection: Supplies M126/N144/L162
  types without merging them with N100..
\item
  \textbf{\texttt{SAMA-D000035}.} Connection: Makes the grammar, type
  system, program custody and exact semantic path explicit..
\item
  \textbf{\texttt{SAMA-D000036}.} Connection: Receives the exact DOS
  definition, coefficient closure and route families..
\item
  \textbf{\texttt{SAMA-D000037}.} Connection: Receives the
  structure-first router and exact fallback policy..
\item
  \textbf{\texttt{SAMA-P000006}.} Connection: Connects N100 computation
  back to the particle, Starbreaker and binding parent route..
\end{itemize}

The immediate handoff to D35 is a language requirement: source identity,
operator type, program tape, state and history must be serialized so the
grammar certificates can select an exact route without hidden UI or
wrapper semantics.

\subsection{Execution architecture and
controls}\label{viii-ch01-execution-architecture-and-controls}

\subsection{Overview}\label{viii-ch01-overview-1}

The central Volume III computational finding is not that a \(2^{100}\)
state space disappeared. It is that the source-derived N100 grammar
carries an exact factorization certificate. The unchanged 100-spin
Hamiltonian has eight conditional components. Evaluating those
components locally and convolving their exact tensors reconstructs the
same complete result as the H14F variable-elimination route, coefficient
for coefficient.

This separates three quantities that are often confused:

\begin{itemize}
\tightlist
\item
  formal state space = every one of 2\^{}100 assignments counted
\item
  requested exact output = the complete retained-role DOS tensor
\item
  fresh local work = 165,920 component assignments
\end{itemize}

The source grammar determines the first two. A certified execution route
changes only the third. H14F remains the exact fallback when a
component, orbit, fiber, separator or rank certificate is absent or
fails an installation gate. H14F's fourteen-worker surface is a frozen
execution contract where declared, not an informal knob that can be
reinterpreted as an arbitrary host setting.

The controls make the result transferable. N105 restores the excluded
source packet; N81 supplies a source-native octahedral comparison; N72
remains a dense contrast; mixed partial-join and connected-ring graphs
test whether the router can preserve exactness when components are
coupled; N96 demonstrates exact cutset branching and fallback. These
controls show what the router is: a typed grammar-to-work compiler, not
a rule that substitutes the N100 graph for other systems.

Historically, H14F and the structure-first receipts belong to the exact
arithmetic/application lineage now preserved in the delegated frozen
\texttt{SLCV21R} base. Frozen \texttt{SLCQ2-RZ} is the sole
source-recorded default SLC selection engine. Q2 consumes the
byte-identical base services; it does not turn H14F into the present
ordinal selector or erase the earlier exact results.

\subsection{Opening question and conceptual
picture}\label{viii-ch01-opening-question-and-conceptual-picture-1}

The governing question is:

\begin{quote}
When a finite grammar defines a complete exact calculation, can the
grammar itself certify a smaller execution plan without changing the
mathematical state space, Hamiltonian or requested output?
\end{quote}

The answer is represented by

\begin{itemize}
\tightlist
\item
  typed source grammar
\item
  declared Hamiltonian and requested tensor
\item
  structural certificate discovery
\item
  admitted exact route
\item
  component convolution
\item
  orbit-sector counting
\item
  separator branching
\item
  rank-aware readout
\item
  or H14F exact fallback
\item
  independent reconstruction of the same tensor
\end{itemize}

The router never asks which route would make a preferred outcome appear.
It asks which exact decomposition is already present in the declared
object. If the certificate is missing, the object is not modified to
create one.

\subsection{Typed definitions and execution
contract}\label{viii-ch01-typed-definitions-and-execution-contract}

\subsubsection{Mathematical object versus execution
profile}\label{viii-ch01-mathematical-object-versus-execution-profile}

For a fixed Hamiltonian \(H\), define the retained-role tensor

\[
G(e,r)=\#\{\sigma:H(\sigma)=e,\ R(\sigma)=r\},
\]

where \(r\) is one of the declared port-role states. Its scalar DOS is

\[
g(e)=\sum_rG(e,r).
\]

The mathematical object consists of \(H\), the assignment domain, the
role map \(R\), and every requested coefficient \(G(e,r)\). H14C and
H14F specify how that object is evaluated. They may change elimination
order, batching, residue storage, worker persistence and cache policy.
They may not change \(H\), delete assignments, reinterpret the roles or
replace \(G\) with a smaller target.

\subsubsection{H14C and H14F}\label{viii-ch01-h14c-and-h14f}

H14C is the completed fourteen-core profile with complete-prime
calibration and the accepted retained-capacity contract. H14F adds
mixed-width compilation, persistent workers, reduced residue storage and
an ephemeral scale-specific projection cache. Both are exact execution
profiles.

Where an H14F campaign declares fourteen pinned physical workers, that
width is part of the frozen architecture. It cannot silently become
``whatever number of host workers is available,'' because batching,
custody and timing receipts are indexed to the declared worker surface.
A new worker mapping would be a successor execution profile, not a
reinterpretation of the old one.

\subsubsection{Structural certificate}\label{viii-ch01-structural-certificate}

A structural certificate is exact evidence that the requested object can
be reconstructed from smaller typed objects. Examples include:

\begin{itemize}
\tightlist
\item
  disconnected components after an explicit conditioning boundary;
\item
  a permutation orbit on which energy depends only on marked population;
\item
  a cutset whose fixed assignments reduce admitted elimination width;
\item
  a fee-exempt fiber whose multiplicity can be lifted exactly; and
\item
  a constrained readout with a known independent rank.
\end{itemize}

The certificate must contain the reconstruction rule. A
low-width-looking plot or favorable runtime is not enough.

\subsubsection{Router and fallback}\label{viii-ch01-router-and-fallback}

The routing policy is

\begin{itemize}
\tightlist
\item
  if conditional components close:
\item
  evaluate each component exactly and convolve
\item
  elif uniform orbit sectors close:
\item
  count exact marked-population sectors
\item
  elif a declared cutset lowers admitted width:
\item
  branch exactly on every cutset assignment
\item
  elif the output is a constrained fiber/readout:
\item
  evaluate independent rank and reconstruct the ambient view
\item
  else:
\item
  execute the exact H14F variable-elimination fallback
\end{itemize}

Every branch returns the complete requested tensor. Failure to certify
one branch is a routing outcome, not permission to approximate.

\subsection{The complete H14F N100
result}\label{viii-ch01-the-complete-h14f-n100-result}

\subsubsection{Source object}\label{viii-ch01-source-object}

The source-derived N100 object has

\begin{itemize}
\tightlist
\item
  \textbf{independent spins.} Declared value: 100.
\item
  \textbf{source edges.} Declared value: 260.
\item
  \textbf{induced width.} Declared value: 5.
\item
  \textbf{retained source packet.} Declared value: \(p=1,g=0\).
\item
  \textbf{role order.} Declared value: plus, minus, anti-plus,
  anti-minus.
\item
  \textbf{complete configurations.} Declared value: \(2^{100}\).
\end{itemize}

The exact configuration count is

\[
2^{100}
=1{,}267{,}650{,}600{,}228{,}229{,}401{,}496{,}703{,}205{,}376.
\]

H14F treats the complete 100-spin, 260-edge Hamiltonian as one
retained-port variable-elimination object. It executes 2,048 singleton
root tasks through four CRT primes.

\subsubsection{Retained-role closure}\label{viii-ch01-retained-role-closure}

Four binary retained roles give \(2^4=16\) port rows. Each row sums to
\(2^{96}\), because fixing the four retained role bits leaves 96
independent binary coordinates. Summing all rows gives

\[
\sum_{r=0}^{15}\sum_eG(e,r)
=16\cdot2^{96}
=2^{100}.
\]

The scalar result has 177 occupied energy bins and ground sector
((-260,256)): energy \(-260\) with degeneracy 256. The conditional port
rows organize into three spectra according to plus-spin count:

\begin{itemize}
\tightlist
\item
  \textbf{0.} Representative state: 0; Ground energy: -260; Ground
  count: 128; Bins: 174.
\item
  \textbf{1.} Representative state: 1; Ground energy: -250; Ground
  count: 128; Bins: 174.
\item
  \textbf{2.} Representative state: 3; Ground energy: -244; Ground
  count: 256; Bins: 171.
\item
  \textbf{3.} Representative state: 7; Ground energy: -250; Ground
  count: 128; Bins: 174.
\item
  \textbf{4.} Representative state: 15; Ground energy: -260; Ground
  count: 128; Bins: 174.
\end{itemize}

Thus counts 0/4 share one spectrum, 1/3 share a second and 2 supplies
the third. All 24 source-role permutations and the global spin
complement close exactly. This is a computationally typed conditional
tensor, not an elementwise F81 physical-port assignment.

\subsubsection{Execution receipt}\label{viii-ch01-execution-receipt}

The four prime-kernel walls were approximately

\[
0.890328, 0.852483, 0.797375, 0.799622\ \text{s},
\]

and the total runner wall was 3.828137395 seconds. The scalar DOS
semantic hash is
\nolinkurl{6e404623175ba60629110fb17f75ca0f0398e690db9de7b7a97766d32d0c2d30}
and the retained tensor semantic hash is
\nolinkurl{33ec2e3f5cd3a1e66eb892730181e89bb1aff71d93b97fbbba4959bb238bd78d}.
These hashes become the equality target for the structure-first route.

\subsection{Deriving the N100 structure-first
work}\label{viii-ch01-deriving-the-n100-structure-first-work}

\subsubsection{Component census}\label{viii-ch01-component-census}

The unchanged conditional grammar separates into eight exact components
of sizes

\[
15,15,15,15,15,10,10,5.
\]

The sizes sum to the original independent rank:

\[
5\cdot15+2\cdot10+5=75+20+5=100.
\]

For a component of \(n_a\) binary variables, direct local enumeration
requires \(2^{n_a}\) assignments. Therefore the total fresh component
work is

\[
W_{100}=5\cdot2^{15}+2\cdot2^{10}+2^5.
\]

Evaluating each term,

\[
5\cdot32{,}768=163{,}840,
\]

\[
2\cdot1{,}024=2{,}048,
\]

and

\[
W_{100}=163{,}840+2{,}048+32=165{,}920.
\]

\subsubsection{\texorpdfstring{Why 165,920 does not replace
\(2^{100}\)}{Why 165,920 does not replace 2\^{}\{100\}}}\label{viii-ch01-why-165920-does-not-replace-2100}

Let component \(a\) return its exact conditional polynomial tensor
\(P_a(z,r_a)\). The full tensor is reconstructed by summing compatible
role assignments and convolving energies:

\[
P_{100}(z,r)
=\sum_{r_1\oplus\cdots\oplus r_8=r}
 \prod_{a=1}^8P_a(z,r_a).
\]

At \(z=1\), each local polynomial counts all local assignments, so

\[
P_{100}(1)
=\prod_{a=1}^82^{n_a}
=2^{\sum_an_a}
=2^{100}.
\]

The 165,920 number counts fresh local assignments evaluated by the
algorithm. The \(2^{100}\) number counts formal global assignments
represented by the reconstructed tensor. They answer different
questions.

\subsubsection{Same-output equality}\label{viii-ch01-same-output-equality}

The structure-first route completed in 0.500951441 seconds. Relative to
the 3.828137395-second H14F route, the observed ratio is

\[
\frac{3.828137395}{0.500951441}\approx7.642.
\]

This comparison is meaningful because the retained-role tensors are
coefficientwise identical and have the same semantic hash. The route did
not change the Hamiltonian, sample states, truncate bins or request only
the ground sector.

\subsubsection{Grammar-to-work
identities}\label{viii-ch01-grammar-to-work-identities}

The work also admits exact source-grammar decompositions:

\[
W_{100}=162\cdot2^{10}+2^5
=2^5(72^2+1).
\]

Indeed,

\[
162\cdot1{,}024+32=165{,}888+32=165{,}920,
\]

and \(72^2=5{,}184=32\cdot162\). Therefore

\[
\frac{W_{100}-2^5}{2^{10}}=162,
\qquad
2^5\cdot162=72^2.
\]

The residual \(2^5\) is the formal assignment space of the surviving
five-role \(p=9,g=1\) packet. It is not X1 and not a physical closure
object. The identity interprets source grammar and workload; it does not
claim that a ledger total was inserted as an energy coefficient.

\subsection{Source-native controls and
deviations}\label{viii-ch01-source-native-controls-and-deviations-1}

\subsubsection{N105 restoration control}\label{viii-ch01-n105-restoration-control}

N100 excludes the historical \(p=9,g=0\) witness packet. Restoring it
produces N105 and changes the exact work to

\[
W_{105}=166{,}912=163\cdot2^{10}.
\]

The difference is

\[
\Delta W=W_{105}-W_{100}
=166{,}912-165{,}920
=992
=2^{10}-2^5.
\]

This is a source-native perturbation: it restores an identified packet
rather than adding an arbitrary edge until timing changes. It also
changes the N100/N81 completion-path capacity from 40 to 45 rows.

\subsubsection{N81 octahedral control}\label{viii-ch01-n81-octahedral-control-1}

The N81 control has component sizes

\[
[11,10,11,10,10,11,10,8]
\]

and local work

\[
\begin{aligned}
W_{81}&=3\cdot2^{11}+4\cdot2^{10}+2^8\\
&=6{,}144+4{,}096+256\\
&=10{,}496=2^7(81+1).
\end{aligned}
\]

Its octahedral bridge retains 48 cube embeddings in one full-octahedral
orbit, the completion path (1-8-9-12), reflected C4 gluing and 49,152
incidences. This is a geometry/compiler control. Its possible causal
relation to runtime remains a research question rather than an installed
explanation.

\subsubsection{N72 dense contrast}\label{viii-ch01-n72-dense-contrast-1}

N72 remains dense. It is not edited to reproduce N100's component
structure. This is the correct contrast because it asks whether the
router recognizes the actual source object rather than imposing one
preferred decomposition on every input.

\subsubsection{Mixed-connectivity
controls}\label{viii-ch01-mixed-connectivity-controls}

The installation tournament adds two controlled deviations:

\begin{itemize}
\tightlist
\item
  a one-edge partial join, width 5, ground sector ((-261,128)); and
\item
  a connected ring, width 5, ground sector ((-268,2)).
\end{itemize}

Both complete sixteen-row tensors equal an independent
original-component bridge contraction coefficient for coefficient. The
partial join retains 128 component-orientation choices; the connected
cyclic skeleton leaves the two global orientations. These tests show
that a changed coupling topology can be handled exactly when its
reconstruction rule is explicit.

\subsubsection{Dense N96 fallback}\label{viii-ch01-dense-n96-fallback}

The full N96 graph remains width 24. A frozen five-site cutset creates

\[
2^5=32
\]

exact branches, each admitted to H14F at width 22. This is neither a
component convolution nor an approximate shortcut. Every cutset
assignment is retained, and H14F performs the exact branch work. The
regression ran 4,096 fresh tasks and retained the frozen N96 numerical
authority.

\subsection{Failure, correction and routing
discipline}\label{viii-ch01-failure-correction-and-routing-discipline}

The principal wrong route would be to treat N100's favorable graph as
the universal computational model. That route is rejected in stages:

\begin{itemize}
\tightlist
\item
  wrong control: substitute the N100 component graph for another grammar
\item
  source identity is lost
\item
  correction: inspect the new object's own
  component/orbit/cutset/fiber/rank certificates
\item
  choose only an admitted exact route
\item
  retest: reconstruct the complete tensor independently
\item
  compare coefficientwise and replay dense regressions
\item
  fallback: if no certificate closes, retain H14F
\end{itemize}

The installation gates execute this logic rather than merely state it.
They record 12/12 route checks, 18/18 mixed arithmetic checks, 9/9
N18/N54 checks, 16/16 N72 checks and 20/20 N96 checks: 75/75 combined
executed checks. Their admission was for a separate owner-authorized
promotion decision; the gate campaign itself did not silently rewrite a
current pointer.

A second deviation is the source-native N100-to-F81 compiler gap. N100
has 100/100 typed source keys and Shared72 has aggregate connector
incidence for 72/100 rows, while CR211 exposes 162 address rows but zero
elementwise QP assignments. The missing elementwise join cannot be
replaced by aggregate counts or external labels.
\texttt{G:G1@SAM-RESEARCH} preserves that partial-domain boundary in
\texttt{SAMA-D000034}; it is not a direct H14F result key for this
chapter.

\subsection{Recorded authority and delegated-base
status}\label{viii-ch01-recorded-authority-and-delegated-base-status}

Historical H14F, structure-first, Exact Write, W9P/Theta, support/Gram,
tensor-orbit and bounded application artifacts remain frozen within the
\texttt{SLCV21R} arithmetic/application lineage. Recorded authority is \texttt{SLCQ2-RZ}, the source-recorded default hybrid ordinal selection engine with 70 exact source-visible features. It delegates byte-identical arithmetic/application services to the frozen \texttt{SLCV21R} base, retaining the exact H14F fallback and structure-first routing lineage.


\texttt{SLCV21R} is noncurrent as a revision, but its declared services
remain the delegated base. Q1 is the immutable fallback. Q2 does not
install physical calibration, Binding, ATOM3D or RH merely by consuming
base arithmetic. H14F therefore remains load-bearing without becoming
the current selection engine.

\subsection{Established result, boundary and forward
handoff}\label{viii-ch01-established-result-boundary-and-forward-handoff}

The complete N100 result counts all \(2^{100}\) assignments of the
unchanged 100-spin, 260-edge Hamiltonian and returns the full exact
retained-role tensor. The structure-first route evaluates 165,920 local
assignments and reconstructs that same tensor coefficientwise. The
execution-work reduction is established for this source grammar; it is
not asserted for every arbitrary 100-spin graph.

The router's general result is procedural and exact: detect the source
certificate, preserve the requested object, reconstruct independently,
and retain H14F when no reduced route closes. N105, N81, N72,
mixed-graph and N96 controls define where each branch applies.

The forward handoff goes to \texttt{SAMA-D000041}, where Starbreaker's
own grammar is routed without substituting the N100 graph. It also
supports the current-Q2 ordinal chapter \texttt{SAMA-D000055}, where
exact base computation and source-visible selection must remain
different types.

